%-------------------------------------------------------------------------------
\subsection{Discussion}
%-------------------------------------------------------------------------------

In this paper, we identify a novel variant of order tampering attacks through formal analysis and demonstrate the attack in real-world payment systems.
Our findings reveal a long-standing anti-pattern in the high-level protocol design of third-party online payment systems, namely the absence of user-side order authentication.
To address this issue, we propose a new payment model that systematically eliminates this design deficiency.

\textbf{Advantages.}
From a protocol design perspective, our payment model exhibits three key advantages.
\begin{enumerate}
  \item \textcolor{red}{\textbf{Minimal trust reduction.}}
At the design level, we reduce reliance on trusted external channels, ensuring payment security even under the weakened external channel assumption.
  \item \textcolor{red}{\textbf{Backward compatibility.}}
Our payment model is highly compatible with existing payment models, preserving their four-phase structure and facilitating incremental deployment and adoption by payment service providers.
  \item \textbf{Formally provable security.} 
Our formal analysis results demonstrate that our payment model effectively prevents order tampering attacks during regular payments and satisfies all security properties defined in this paper.
\end{enumerate}

\textbf{Limitations.} 
We honestly present the limitations and trade-offs of our payment model.
\begin{enumerate}
  \item \textbf{Signature and storage overhead.}
Our payment model introduces one additional signature operation, and requires both merchants and users to locally store extra keys, which incurs additional computational and implementation overhead.
However, no new cryptographic primitives are introduced, and we consider the resulting overhead to be acceptable in practice.
  \item \textbf{Additional user interaction during the initial payment.}
During the initial payment, users are required to explicitly confirm their trust in the merchant, introducing additional user interaction.
For regular payments, no extra interaction is required, and the overall user experience remains comparable to that of existing payment models.
  \item \textbf{No security improvement for the initial payment.}
Our proposed model does not provide additional security guarantees for the initial payment.
Instead, our design follows the widely adopted Trust-On-First-Use (TOFU) model, under which the initial interaction may be vulnerable to active network attacks.
Once a legitimate payment context is established, however, any subsequent deviation can be reliably detected.
Consequently, while an adversary may interfere with the initial payment, our scheme guarantees the integrity of subsequent payments and enables attack detection even if the first payment was compromised.
\end{enumerate}